%==============================================================================
% Книга: Как устроен наш общий дом
% Путешествие по Вселенной без тайн и невидимых монстров
% Автор: Виктор Логвинович
%==============================================================================
\documentclass[12pt,a4paper,oneside]{book}

% Кодировки и язык
\usepackage[T2A]{fontenc}
\usepackage[utf8]{inputenc}
\usepackage[russian]{babel}
\usepackage{indentfirst}

% Оформление
\usepackage{geometry}
\geometry{margin=2.5cm, top=3cm, bottom=3cm}
\setlength{\headheight}{28pt} % Увеличена высота для длинных заголовков
\usepackage{microtype}
\usepackage{titlesec}
\usepackage{fancyhdr}
\usepackage{amsmath, amssymb}
\usepackage{array}
\usepackage{booktabs}
\usepackage{graphicx}
\usepackage{longtable} % Пакет для разрыва длинных таблиц на страницы

% Гиперссылки (загружаем последними)
\usepackage{hyperref}
\hypersetup{colorlinks=true, linkcolor=black, urlcolor=blue}

% Настройка заголовков
\titleformat{\chapter}[display]
  {\normalfont\huge\bfseries\centering}{\chaptertitlename\ \thechapter}{20pt}{\Huge}
\titlespacing*{\chapter}{0pt}{40pt}{30pt}
\titleformat{\section}
  {\normalfont\Large\bfseries}{\thesection}{1em}{}
\titleformat{\subsection}
  {\normalfont\large\bfseries}{\thesubsection}{1em}{}

% Колонтитулы
\pagestyle{fancy}
\fancyhf{}
\fancyhead[LE,RO]{\thepage}
\fancyhead[RE]{\leftmark}
\fancyhead[LO]{\rightmark}
\fancyfoot[C]{\small Парадигма IT3}

%==============================================================================
\begin{document}

%==============================================================================
\begin{titlepage}
    \centering
    \vspace*{2cm}
    {\Huge\bfseries Как устроен наш общий дом\par}
    \vspace{0.5cm}
    {\Large Путешествие по Вселенной без тайн и невидимых монстров\par}
    \vspace{2cm}
    {\Large Виктор Логвинович\par}
    \vfill
    {\large Научно-популярное издание на основе серии IT3\par}
    \vspace{1cm}
    {2026\par}
\end{titlepage}

%==============================================================================
\frontmatter
\tableofcontents

%==============================================================================
\chapter*{Введение: Пазл, в котором не хватает кусочков}
\addcontentsline{toc}{chapter}{Введение: Пазл, в котором не хватает кусочков}

Представьте, что вы собираете гигантский пазл с картинкой ночного неба. Почти всё на месте, но в центре не хватает трёх фрагментов. Вы ищете их под диваном, в коробке, на полу — нигде нет. И тогда вы говорите: «Ладно. Значит, этих кусочков просто не существует. Но чтобы пазл не развалился, мы добавим невидимый клей и невидимые распорки».

Звучит странно, правда? Но именно так сегодня выглядит главная космологическая модель. Астрономы видят, что галактики вращаются быстрее, чем должны. Вместо того чтобы пересмотреть, как устроено само пространство, мы говорим: «Там есть невидимая \textbf{Тёмная Материя}». Мы видим, что Вселенная расширяется с ускорением. Вместо поиска причины в форме космоса, мы говорим: «Там действует невидимая \textbf{Тёмная Энергия}».

Мы усложняем мир, населяя его невидимыми сущностями, только ради того, чтобы спасти старую привычку думать, будто Вселенная бесконечна во все стороны.

Эта книга предлагает другой путь. Путь простоты.
Что если убрать весь этот «тёмный» мусор? Что если мы просто неправильно представили себе форму нашего космического дома?

Добро пожаловать в парадигму IT3. Здесь нет магии. Здесь есть только геометрия, логика и удивительная гармония конечного мира. Поехали.

%==============================================================================
\mainmatter

%==============================================================================
\chapter{Форма нашего дома: Конечность без краёв}

\section{Бесконечность или комната без стен?}

Мы привыкли думать, что Вселенная бесконечна. Это удобно. Если космос не имеет границ, свет может лететь вечно, ни во что не упираясь. Гравитация не должна «замыкаться».

Но у бесконечности есть цена. Если Вселенная бесконечна, в ней должны существовать волны абсолютно любого размера. Представьте океан. В бесконечном океане могут гулять цунами размером с континент. Однако, когда астрономы смотрят на реликтовое излучение (самое древнее «эхо» Большого взрыва), они видят странную вещь: самых крупных температурных волн там нет. Их мощность на 20\% ниже, чем предсказывает теория.

Почему? Потому что волны не помещаются в «комнату».

Если вы натянете гитарную струну и дёрнете её, она не сможет издать волну длиннее самой себя. Самые низкие басы просто не влезут в короткую струну. Вселенная работает точно так же. В ней не могут существовать флуктуации (колебания), которые больше её собственного размера. И данные спутника Planck показывают именно это: обрезка на самых больших масштабах. У Вселенной есть конкретный размер.

Когда учёные подставляют эти данные в уравнения, математика строго ограничивает этот размер. Получается число: $L_x = 28{,}57$ миллиардов световых лет. Это не случайная цифра. Это ровно тот масштаб, который нужен, чтобы «отрезать» лишние низкие частоты в реликтовом излучении, но при этом не сломать всё остальное.

\section{Эффект «Пакмана»}

Но если Вселенная конечна, то где её край? Что будет, если долететь до стены?

Никакой стены нет. Вспомните старую аркадную игру. Если ваш герой убегает за правый край экрана, он тут же появляется с левой стороны. Если улетает вверх — выныривает снизу. Карта замкнута, но у неё нет границ. Математики называют такую форму \textbf{тором} (или бубликом).

\textbf{Что такое тор простым языком?} Это замкнутое пространство, у которого нет краёв. Вы можете лететь вперёд вечно и в итоге вернётесь в точку старта. Наша Вселенная — это трёхмерный тор.

Это меняет всё. В бесконечной Вселенной энергия и гравитация рассеиваются в никуда. В конечном торе они не могут исчезнуть. Они отражаются, накладываются сами на себя и заставляют пространство «гудеть», как стенки музыкального инструмента.

\section{Почему \texorpdfstring{$\sqrt{2}$ и $\sqrt{3}$}{корень из 2 и корень из 3}?}

Здесь вступает в игру вторая деталь. Если стороны нашего «космического ящика» относятся как простые числа (например, $1:2:3$), свет будет отражаться по одним и тем же траекториям, создавая искусственные узоры и эхо. Чтобы Вселенная выглядела случайной и однородной, стороны должны относиться друг к другу как \textbf{иррациональные числа}: $1 : \sqrt{2} : \sqrt{3}$.

\textbf{Что такое иррациональные соотношения?} Это пропорции, которые нельзя выразить простой дробью. Они гарантируют, что траектории света и гравитационных волн никогда не замкнутся идеально. Они заполняют пространство равномерно, имитируя наблюдаемую случайность, сохраняя при этом конечность. Никакой мистики. Чистая математика эргодичности: детерминированный хаос, который выглядит как идеальный беспорядок.

%==============================================================================
\chapter{Космическое эхо и пропавшие басы}

\section{Реликтовое излучение: фотография младенческой Вселенной}

Когда Вселенной было всего 380 000 лет, она остыла достаточно, чтобы свет смог свободно лететь сквозь пространство. Этот древний свет до сих пор путешествует вокруг нас. Мы называем его \textbf{реликтовым излучением} (или CMB). Это как Baby-фотография космоса.

На этой фотографии видны крошечные пятна: где-то горячее, где-то холоднее. Учёные раскладывают их на «ноты» (мультиполи). Низкие ноты (низкие $\ell$) — это самые длинные, плавные волны. Высокие ноты — короткие, частые ряби.

\section{Аномалия низких мультиполей}

Стандартная модель предсказывает, сколько энергии должно быть в каждой «ноте». Но самые низкие, самые глубокие «басы» оказались тише, чем ожидалось. Их не хватает.

В бесконечной Вселенной это загадка. В нашей конечной — это естественный факт. \textbf{Инфракрасный обрез} (научный термин для «отсечения длинных волн») работает как фильтр: если волна не помещается в размер тора, она просто не возникает. Музыкальный ящик не может издать звук ниже своей собственной длины. Вселенная подтверждает это наглядно.

%==============================================================================
\chapter{Парадокс спидометра: Почему правы обе стороны?}

\section{Напряжение Хаббла}

Одна из самых горячих загадок современной физики — разброс в измерениях скорости расширения Вселенной. Эта скорость называется \textbf{постоянной Хаббла} ($H_0$).

\begin{itemize}
    \item Если смотреть на ранний космос (реликтовое излучение), получается около \textbf{67 км/с на мегапарсек}.
    \item Если смотреть на близкие сверхновые и цефеиды (звёзды-маяки), получается около \textbf{73}.
\end{itemize}

Кто-то ошибается? Или данные врут? Это расхождение называют \textbf{напряжением Хаббла} (Hubble tension).

\section{Решение через «фазы» дома}

В парадигме IT3 \textbf{правы обе стороны}. Они просто измеряют Вселенную в разных фазах её взросления.

Ранние зонды видят космос до того, как сформировалась крупная структура. Тогда пространство было гладким. Поздние зонды видят космос после того, как образовались гигантские пустоты и нити галактик. В этот момент включается особый механизм, который ускоряет расширение.

Это не конфликт данных. Это подпись эволюции структуры. Одна Вселенная, два режима работы. Как если бы вы измеряли скорость реки в сухом русле и во время паводка: показания разные, но река одна.

%==============================================================================
\chapter{Космическая губка: Когда пустота начинает работать}

\section{Войды: больше, чем просто пустота}

Когда астрономы картографируют Вселенную, они видят не равномерный туман из звёзд, а гигантскую паутину. Галактики собраны в светящиеся нити, а между ними зияют невероятные пустоты — \textbf{войды}. Они занимают около \textbf{75\% объёма космоса}.

Раньше считалось, что эти пустоты — просто «шум», побочный продукт гравитации. Но представьте себе обычную кухонную губку. Она состоит из твёрдого каркаса и дырочек. Если налить на неё воду, она потечёт не где попало, а по каналам с наименьшим сопротивлением.

Так работает и космос. Войды — это не отсутствие чего-то. Это \textbf{каналы}, по которым энергия течёт совершенно иначе. Вместе они образуют \textbf{Космическую Губку}.

\section{Как Вселенная нажала на газ (Перколяция)}

В самом начале времён войдов не было. Материя была распределена почти идеально ровно. Губка была «выключена». Но по мере расширения гравитация собирала галактики в кучи, а пустоты росли. Примерно 7–8 миллиардов лет назад войды соединились в единую глобальную сеть. Произошёл фазовый переход, который физики называют \textbf{перколяцией}.

\textbf{Что такое перколяция?} Это момент, когда разрозненные пустоты вдруг соединяются в сквозные пути. Как вода, которая наконец находит путь сквозь кофейную гущу в турке. Или как пена в ванной, которая вдруг становится единой массой.

В этот момент «космическая губка» включилась. Она начала усиливать топологические эффекты почти в четыре раза. И именно это включение меняет историю расширения Вселенной в позднюю эпоху.

\section{Топологический сток: Замена тёмной энергии}

Когда структуры формируются, они выделяют гравитационную энергию. В бесконечном пространстве она улетает в вечность. В конечном торе она не может уйти. Она отражается, интерферирует и фокусируется в симметричных точках пространства. Это создаёт макроскопическое \textbf{отрицательное давление}. Мы называем это \textbf{топологическим стоком}.

Простыми словами: материя не просто разбавляется при расширении, часть её энергии непрерывно «перетекает» в геометрию пространства, создавая эффект отталкивания. Когда губка включилась, этот сток усилился. Пространство начало расширяться быстрее не потому, что в него влили новую «тёмную» энергию, а потому что изменилась проводимость самой ткани реальности. Ускорение без магии. Только конечность и губчатая структура.

%==============================================================================
\chapter{Космическое трение: Почему галактики не слипаются?}

\section{Напряжение S8}

Есть ещё одна загадка: \textbf{напряжение S8}. Компьютерные симуляции говорят, что галактики должны собираться в скопления активнее, чем мы видим на самом деле. Реальность показывает меньше «слипания».

\section{Геометрическое трение}

В нашей модели тот же самый топологический сток добавляет дополнительное «трение» в рост структур. Представьте, что вы пытаетесь скатать снежок, но снег немного влажный и липкий: ком растёт медленнее. Так и здесь: часть энергии, которая должна была идти на рост галактик, уходит в геометрию тора через губку.

Это подавляет формирование структур примерно на 10\%. Результат? Теория идеально совпадает с наблюдениями \textbf{слабого гравитационного линзирования} (когда свет далёких галактик слегка искажается, как в кривом зеркале, позволяя нам «взвесить» невидимую массу). Тёмной материи как частицы не нужно. Есть только кривизна путей в конечном резонаторе.

%==============================================================================
\chapter{Музыка частиц: Откуда берётся масса?}

\section{Бозон Хиггса и его 125 ГэВ}

В физике элементарных частиц есть знаменитая загадка: почему бозон Хиггса весит ровно 125 ГэВ? В стандартной модели это число просто вписано от руки. Оно кажется случайным.

Но вспомним, что наша конечная Вселенная похожа на музыкальный инструмент. Если вы возьмёте трубу разной длины, она будет звучать на разных нотах. Вакуум в космосе не пустой. Он наполнен квантовыми флуктуациями, которые тоже подчиняются правилам резонатора. Их энергия, а значит и массы частиц, которые с ними взаимодействуют, зависят от размера «ящика».

\section{Эмерджентность: Свойство, которое рождается вместе}

Оказывается, масса Хиггса не «вшита» в частицу изначально. Она \textbf{эмерджентна}.

\textbf{Что такое эмерджентность?} Это свойство, которое появляется только когда много частей работают вместе, как волна на стадионе: один человек не волна, но тысяча людей — уже волна. Так и масса: она возникает как коллективный эффект взаимодействия топологии пространства, перколяционной сети войдов и голографического предела вакуума.

Математика показывает простую связь: если подставить размер тора ($28{,}57$ Гпк), фактор усиления губки ($\approx 3{,}8$) и геометрические коэффициенты иррационального тора, получается:
\[
m_H \approx 125{,}1 \text{ ГэВ}.
\]
Совпадение с экспериментом не случайное. Это значит, что масса Хиггса — это «нота», которую издаёт наша Вселенная благодаря своему размеру и форме.

%==============================================================================
\chapter{Как мы проверяем: Наука не верит на слово}

Хорошая теория — не та, которую нельзя опровергнуть, а та, которую \textit{можно} проверить. Это называется \textbf{фальсифицируемость}. Парадигма IT3 не прячется за красивыми словами. Она делает чёткие предсказания:

\begin{enumerate}
    \item \textbf{Совпадающие окружности в CMB.} Если Вселенная замкнута, свет от одних и тех же областей должен приходить к нам по разным маршрутам. На карте реликтового излучения должны остаться пары кругов с одинаковой температурной структурой. IT3 предсказывает их угловой радиус: $30^\circ$–$60^\circ$. Будущие телескопы (CMB-S4, LiteBIRD) будут искать именно их.
    
    \item \textbf{Эволюция постоянной Хаббла.} Если губка включается постепенно, то выводимая $H_0$ должна плавно расти при красных смещениях $z < 1$. Обзоры DESI и Euclid измерят это с точностью до процента.
    
    \item \textbf{Подавление гравитационных волн.} Инфляция предсказывает заметный фон первичных волн. В конечном торе длинные волны просто не помещаются. IT3 предсказывает, что их почти не должно быть ($r < 0.01$).
    
    \item \textbf{Аномалии в сверхвойдах.} Поток энергии через перколирующие пустоты должен оставлять след в реликтовом излучении. Его амплитуда должна быть примерно вдвое больше, чем в старой модели.
\end{enumerate}

Если эти сигнатуры не будут найдены, парадигма будет уточнена или заменена. Так и работает наука. Но пока звёзды молчат, уравнения уже сложились в самосогласованную картину.

%==============================================================================
\backmatter

%==============================================================================
\chapter*{Заключение: Один дом, одна геометрия}
\addcontentsline{toc}{chapter}{Заключение: Один дом, одна геометрия}

Современная космология десятилетиями добавляла невидимые компоненты, чтобы сохранить допущение: пространство бесконечно. Парадигма IT3 задаёт более простой вопрос: \emph{что если пространство конечно, и мы просто забыли учесть, как энергия течёт в замкнутой форме?}

Ответ удивительно согласован:
\begin{itemize}
    \item \textbf{Форма} отсекает самые длинные волны, объясняя аномалию низких мультиполей.
    \item \textbf{Поток} отводит энергию в топологические пучности, заменяя тёмную энергию физическим механизмом диссипации.
    \item \textbf{Структура} каналирует этот поток через войды, активируя ускорение и подавляя избыточное кластеризование.
    \item \textbf{Резонанс} вакуума в конечном торе задаёт массу фундаментальных частиц.
\end{itemize}

Никаких новых частиц. Никакой тонкой настройки. Никакой мультивселенной. Только геометрия конечного тора, термодинамика диссипации и наблюдаемое губчатое распределение материи. Один параметр ($L_x$) вместо десятка подгонок.

Следующее поколение обзоров решительно проверит эти предсказания. Если данные совпадут, космология перейдёт от парадигмы невидимых субстанций к парадигме видимой геометрии. Если нет — модель будет уточнена. Это не поражение. Это прогресс.

Но на данный момент теоретический аппарат завершён, самосогласован и готов к тому, чтобы природа вынесла свой вердикт. Вселенная может быть не такой загадочной, как нам казалось. Возможно, ей просто нужна была правильная форма, чтобы всё встало на свои места.

\vspace{2cm}
\begin{flushright}
\textit{«Пока учёные ищут окружности, геометрия уже работает.»}
\end{flushright}

%==============================================================================
\chapter*{Глоссарий для юных исследователей}
\addcontentsline{toc}{chapter}{Глоссарий для юных исследователей}

Краткий справочник всех научных терминов из книги.

\renewcommand{\arraystretch}{1.3}
\begin{longtable}{>{\bfseries}p{3.5cm}p{11cm}}
\toprule
\textbf{Термин} & \textbf{Простое объяснение} \\
\midrule
\endfirsthead

\multicolumn{2}{c}{{\small\textit{Продолжение глоссария}}}\\
\toprule
\textbf{Термин} & \textbf{Простое объяснение} \\
\midrule
\endhead

\bottomrule
\endfoot

\bottomrule
\endlastfoot

\textbf{Космология} & Наука о Вселенной в целом: как она возникла, как устроена и как будет развиваться. \\
\textbf{Топология} & Раздел математики, изучающий, как пространство соединено само с собой (например, замкнуто ли оно, есть ли у него «дырки» или края). \\
\textbf{Тор (3-тор)} & Замкнутое пространство без краёв. Как экран старой видеоигры: уходишь за правый край — появляешься слева. \\
\textbf{Иррациональные соотношения} & Пропорции сторон пространства, которые нельзя выразить простой дробью (например, $1:\sqrt{2}:\sqrt{3}$). Они гарантируют, что свет и волны никогда не зациклятся одинаково, создавая видимую случайность. \\
\textbf{Реликтовое излучение (CMB)} & Древнейший свет во Вселенной, «эхо» Большого взрыва. Доходит до нас отовсюду и показывает, какой была Вселенная в младенчестве. \\
\textbf{Низкие мультиполи (низкие $\ell$)} & Самые длинные, плавные волны температуры в реликтовом излучении. Аналоги «басовых нот» в космической симфонии. \\
\textbf{Инфракрасный обрез} & Естественное ограничение: в конечном пространстве не могут существовать волны длиннее самого пространства. Длинные «басы» просто не помещаются. \\
\textbf{Постоянная Хаббла ($H_0$)} & Скорость, с которой расширяется Вселенная прямо сейчас. Измеряется в километрах в секунду на каждый мегапарсек расстояния. \\
\textbf{Напряжение Хаббла} & Разногласие между двумя способами измерения скорости расширения: по раннему космосу ($\sim 67$) и по близким объектам ($\sim 73$). В IT3 это не ошибка, а признак смены «фазы» Вселенной. \\
\textbf{Войды} & Гигантские космические пустоты, где почти нет галактик. Занимают $\sim 75\%$ объёма Вселенной. \\
\textbf{Космическая губка} & Крупномасштабная структура Вселенной, где галактики образуют нити, а войды — пустоты. Работает как волновод, усиливая потоки энергии. \\
\textbf{Перколяция} & Момент, когда разрозненные пустоты соединяются в единую сеть. Как вода, находящая сквозной путь в кофе. В космосе это «включает» ускоренное расширение. \\
\textbf{Топологический сток} & Механизм, при котором энергия в замкнутом пространстве не рассеивается в бесконечность, а фокусируется самой геометрией, создавая отрицательное давление и отталкивание. Заменяет «тёмную энергию». \\
\textbf{Тёмная энергия / Тёмная материя} & Гипотетические невидимые компоненты, которые в старой модели добавляли «вручную», чтобы объяснить ускорение и вращение галактик. В IT3 их роль играет геометрия и структура пространства. \\
\textbf{Напряжение S8} & Наблюдаемое несоответствие: галактик в скоплениях меньше, чем предсказывают симуляции. В IT3 объясняется «космическим трением» от топологического стока. \\
\textbf{Слабое гравитационное линзирование} & Искривление света далёких галактик под действием гравитации. Позволяет «взвешивать» невидимую массу по тому, как она гнёт свет. \\
\textbf{Бозон Хиггса} & Элементарная частица, связанная с полем, которое придаёт массу другим частицам. \\
\textbf{Эмерджентность} & Свойство системы, которое возникает только при взаимодействии многих частей (как температура возникает от движения молекул). Масса в IT3 не «вшита» в частицы, а рождается из резонанса пространства. \\
\textbf{Фальсифицируемость} & Свойство научной теории быть проверяемой и потенциально опровержимой наблюдениями. Если предсказания не сбудутся — теорию изменят или отбросят. \\
\textbf{Совпадающие окружности} & Ожидаемые узоры в реликтовом излучении: если Вселенная замкнута, свет обходит её разными путями, и на небе должны появиться парные круги с одинаковой температурой. \\
\end{longtable}

\vspace{2cm}
\begin{flushright}
\small Апрель 2026 г. | Кобрин, Беларусь\\
\small Репозиторий с кодом и данными: \url{https://github.com/Viktar-Pi/FlatIrrationalTorus}
\end{flushright}

%==============================================================================
\end{document}